%=============================================================================
% FairTox -- contribution report (required in the repository root)
%
% Build:  latexmk -pdf contribution_report.tex
%
% The brief asks for one table with Member Name / Code / Report Sections /
% Experiments / Slides / Other, plus a short description per member, agreed by
% all members. The instructor cross-checks it against the Git history, so the
% module ownership here is the same ownership the pull-request sequence follows
% and the same one the defence deck is split along.
%=============================================================================
\documentclass[10pt]{article}

\usepackage[a4paper,landscape,margin=1.5cm]{geometry}
\usepackage[T1]{fontenc}
\usepackage[utf8]{inputenc}
\usepackage{helvet}
\renewcommand{\familydefault}{\sfdefault}
\usepackage{booktabs}
\usepackage{array}
\usepackage{xcolor}
\usepackage[hidelinks]{hyperref}

\definecolor{fairblue}{HTML}{3B5BA5}
\definecolor{fairgrey}{HTML}{6B7280}

\setlength{\parindent}{0pt}
\pagestyle{empty}

\newcommand{\person}[2]{\textbf{#1}\\[0.2ex]{\footnotesize\color{fairgrey}#2}}

\begin{document}

{\Large\bfseries\color{fairblue} Contribution report}\\[0.6ex]
{\normalsize FairTox --- \emph{Parity, Utility and Safety Are Not Independent}}\\[0.3ex]
{\footnotesize\color{fairgrey}
DLE-AI-202, Cohort I 2026 \,·\, Track 2 \,·\,
\url{https://github.com/nigarrustamova/fairtox} (tag \texttt{v1.0-final})}

\vspace{0.8em}

\footnotesize
\renewcommand{\arraystretch}{1.15}
\setlength{\tabcolsep}{4pt}
\begin{tabular}{%
  >{\raggedright\arraybackslash}p{0.115\textwidth}
  >{\raggedright\arraybackslash}p{0.245\textwidth}
  >{\raggedright\arraybackslash}p{0.165\textwidth}
  >{\raggedright\arraybackslash}p{0.185\textwidth}
  >{\raggedright\arraybackslash}p{0.105\textwidth}
  >{\raggedright\arraybackslash}p{0.115\textwidth}}
\toprule
\textbf{Member} & \textbf{Code} & \textbf{Report sections} &
\textbf{Experiments} & \textbf{Slides} & \textbf{Other} \\
\midrule

\person{Aliyev, Orkhan}{\texttt{@mrorxan}} &
Stage registry and scaffolding; loss weighting and Group DRO
(\texttt{training/losses.py}); comparability and reuse fingerprints
(\texttt{provenance.py}); benchmark bias metrics
(\texttt{fairness/jigsaw\_metrics.py}); controlled comparison
(\texttt{stages/compare.py}); ablation (\texttt{stages/ablate.py}); mechanism
analysis (\texttt{stages/mechanism.py}); cross-family statistics
(\texttt{evaluation/aggregate.py}, \texttt{scripts/summarise\_study.py});
run orchestration (\texttt{pipeline.py}) &
Method; Results A--D, F--H (headline pair, all eight configurations, mechanism,
threshold matching, Group DRO, held-out axes, ablation); Discussion;
Conclusion &
Headline pair over eight seeds; fixed-epoch control; Group DRO; held-out-axis
arm; $\alpha$ sweep and the \texttt{class} control; mechanism and
matched-threshold analyses &
15--19 (mechanism, the loss, headline, all families, the three controls) &
Experiment plan and pre-registered outcomes; scheduling of the booked GPU
window \\
\addlinespace

\person{Mirzayeva, Laman}{\texttt{@lmnmrzyv}} &
Global classification metrics (\texttt{evaluation/metrics.py}); stratified
split and fingerprint (\texttt{data/split.py}); subgroup auditor with Wilson
intervals and support floors (\texttt{fairness/auditor.py}); all plotting
(\texttt{utils/plotting.py}); audit stage; error analysis and the coded sample;
per-subgroup thresholds (\texttt{stages/postproc.py}); duplicate check
(\texttt{stages/leakage.py}); optional analyses; CI workflow &
Problem Formulation (two errors, three axes, the three guards,
pre-registration); Results I--M (threshold-free metrics, post-processing, the
manual coding, who is affected, integrity and leakage); Limitations &
Subgroup audits and threshold sensitivity; the training-free post-processing
arm; the leakage measurement; manual coding of 60 identity-bearing false
positives &
7--10 (three axes, five guards, what the errors are, two coded examples) &
Figure review before submission; test coverage for two other members' modules
\\
\addlinespace

\person{Rustamova, Nigar}{\texttt{@nigarrustamova}} &
Repository, branch protection and release tag; logging
(\texttt{utils/logging\_utils.py}); configuration system
(\texttt{config.py}, \texttt{configs/default.yaml}, \texttt{smoke.yaml});
artefact I/O and prediction format (\texttt{utils/io.py},
\texttt{evaluation/predictions.py}); corpus loader and synthetic corpus
(\texttt{data/load.py}); EDA and TF-IDF baseline stages; command-line entry
point (\texttt{scripts/run\_all.py}, \texttt{run\_all.sh}); \texttt{README.md}
&
Introduction; Related Work (measuring unintended bias, and the gap this paper
fills); Experimental Setup (corpus and splits, implementation and
reproducibility); Tools and Acknowledgements &
Corpus EDA and subgroup support counts; the TF-IDF + logistic-regression
baseline; the synthetic smoke pipeline used to validate the code path &
3--6 (the problem, the three intervention points, the corpus, the classical
baseline) &
Repository owner and submitter; one-command reproduction path; documentation
review \\
\addlinespace

\person{Samadov, Nijat}{\texttt{@NijatSamadov}} &
Device and seed utilities; tokenisation and dataset
(\texttt{data/dataset.py}); classifier head (\texttt{models/transformer.py});
fine-tuning loop (\texttt{training/trainer.py}); training stage
(\texttt{stages/train.py}); counterfactual augmentation
(\texttt{data/augment.py}); token attribution
(\texttt{explain/attributions.py}); the eight booked-window configuration
files; \texttt{TRAINING\_DAY.md} and the Kaggle runner notebook &
Related Work (pre-processing); Experimental Setup (model and training,
compute); Results E (counterfactual augmentation, dose response) &
All 115 trainings inside the booked window; the DistilRoBERTa backbone arm;
counterfactual swapping at both doses; the five-run noise-floor measurement &
11--14 (the model, 115 models and the noise floor, the swap, the null result)
&
Server environment (MIG targeting, CUDA build); recovery of the interrupted
window \\
\bottomrule
\end{tabular}

\newpage

{\normalsize\bfseries\color{fairblue} Short description per member}

\vspace{0.5em}
\footnotesize
\begin{tabular}{%
  >{\raggedright\arraybackslash}p{0.14\textwidth}
  >{\raggedright\arraybackslash}p{0.82\textwidth}}
\textbf{Aliyev, Orkhan} &
Owned the objective and the comparison machinery: the per-sample weighting, the
Group DRO control, the fingerprints that make a controlled pair legitimate, and
the stages that turn two arms into a verdict. Wrote the mechanism analysis --- the
regression of subgroup FPR on subgroup toxic rate --- and the cross-family
aggregation that produces every table in the paper. Ran the eight-seed headline
replication and the three controls behind the paper's central claim. \\
\addlinespace
\textbf{Mirzayeva, Laman} &
Owned measurement. Wrote the metrics, the stratified split, and the subgroup
auditor, including the two decisions that keep the fairness numbers honest:
separate support floors for FPR and FNR, and Wilson intervals rather than a
bootstrap. Wrote every figure routine, the error-analysis stage, and the
training-free post-processing arm, and read and coded all 60 sampled false
positives by hand. \\
\addlinespace
\textbf{Rustamova, Nigar} &
Set up the repository and the infrastructure the rest of the work sits on:
configuration, logging, artefact I/O, the prediction format, the corpus loader
and the driver that makes the study one command from a clean checkout. Wrote
the EDA and the classical TF-IDF baseline the deep model is measured against,
and the README. Submits the deliverables on behalf of the team. \\
\addlinespace
\textbf{Samadov, Nijat} &
Owned the model and the training path: tokenisation, the classifier head, the
fine-tuning loop and the training stage, plus the configuration files for the
booked window. Ran the window itself --- 38 runs, 115 models, 34.4 GPU-hours ---
and diagnosed the CUDA and MIG problems that blocked it. Wrote the
counterfactual-augmentation arm and the token-attribution analysis. \\
\end{tabular}

\vspace{1.6em}

{\footnotesize
\textbf{Balance and cross-review.} No directory is owned by one person:
\texttt{stages/} and \texttt{tests/} are shared by all four, and
\texttt{utils/}, \texttt{data/}, \texttt{evaluation/}, \texttt{fairness/} and
\texttt{training/} each have two or three owners. Tests are deliberately written
by someone other than the module's author in eight pull requests, so every
member has read and tested at least two other members' code. Every member
reviewed pull requests from every other member.

\vspace{0.6em}
\textbf{Agreement.} All four members have read this report and agree that it
describes the division of work accurately. Prepared for the final submission of
DLE-AI-202, September 2026.}

\end{document}
